\frontsection{RESUMEN}
La anotación de videos en Lengua de Señas Ecuatoriana (LSEc) exige herramientas que respeten el criterio lingüístico del investigador y, al mismo tiempo, permitan aprovechar modelos modernos de Inteligencia Artificial. ELAN, una de las herramientas más utilizadas para la anotación multimodal, facilita segmentar y etiquetar eventos sobre una línea de tiempo. No obstante, su arquitectura basada en Java y su mecanismo de extensión dificultan la conexión directa con servicios de inferencia desarrollados en entornos actuales, principalmente Python, con dependencias especializadas. Esto limita la incorporación de procesos automáticos, como la segmentación temporal y la clasificación de glosas, al flujo habitual de anotación.

En este TIC se desarrolló un middleware de orquestación local para conectar ELAN con servicios de Inteligencia Artificial sin alterar profundamente su arquitectura base. El componente implementado expone una API REST construida con FastAPI y Pydantic, valida solicitudes y manifiestos, mantiene un registro local de modelos instalados y coordina backends ejecutados en contenedores Docker. Además, se integró en ELAN un reconocedor Java basado en AVATecH, encargado de consultar los modelos disponibles, enviar solicitudes de inferencia al middleware y transformar las respuestas en segmentos temporales editables por el investigador.

La metodología siguió un enfoque incremental e iterativo. Primero se revisó la documentación y el código fuente de ELAN para identificar los contratos de extensión, el ciclo de vida de los reconocedores y el formato esperado de los resultados. A partir de ese análisis se diseñó una arquitectura de tres capas: ELAN como cliente de anotación, el middleware como orquestador local y los backends como servicios de inferencia contenerizados. La implementación se organizó por fases, desde la base HTTP y la instalación de paquetes ZIP (modelos), hasta la ejecución con Docker, la cola de inferencia, la observabilidad, la integración gráfica en ELAN y la validación funcional.

Los resultados confirmaron la viabilidad técnica de la integración. El middleware permitió instalar modelos, verificar su disponibilidad, ejecutar inferencias, reportar estados, diferenciar errores funcionales de timeouts y conservar métricas de diagnóstico. Las pruebas de caja negra y caja blanca evidenciaron un comportamiento estable ante modelos desactivados, videos inexistentes, contenedores detenidos, reinstalaciones duplicadas y respuestas inválidas. La principal contribución es una arquitectura local, desacoplada y extensible que incorpora modelos de IA al flujo de ELAN, preserva la soberanía de los datos y separa con claridad las responsabilidades de anotación lingüística, orquestación e inferencia.

\textbf{PALABRAS CLAVE:} ELAN, middleware, orquestación de servicios, Inteligencia Artificial, Lengua de Señas Ecuatoriana.

\frontsection{ABSTRACT}
Video annotation in Ecuadorian Sign Language (LSEc) requires tools that preserve the researcher's linguistic judgment while making it possible to use modern Artificial Intelligence models. ELAN is widely used for multimodal annotation because it supports segmenting and labeling events on a timeline. However, its Java-based architecture and extension mechanism make direct integration difficult with inference services developed in current environments, especially Python, with specialized dependencies. This limits bringing automatic processes, such as temporal segmentation and gloss classification, into the regular annotation workflow.

This Curricular Integration Work developed a local orchestration middleware to connect ELAN with Artificial Intelligence services without deeply modifying the annotation tool. The implemented component exposes a REST API built with FastAPI and Pydantic, validates requests and manifests, maintains a local registry of installed models, and coordinates backends running in Docker containers. Additionally, a Java recognizer based on AVATecH was integrated into ELAN to retrieve available models, send inference requests to the middleware, and transform the responses into editable temporal segments.

The methodology followed an incremental and iterative approach. First, ELAN documentation and source code were analyzed to identify extension contracts, the recognizer life cycle, and the expected result format. Based on this analysis, a three-layer architecture was designed: ELAN as the annotation client, the middleware as the local orchestrator, and model backends as containerized inference services. The implementation was organized into phases covering the HTTP foundation, ZIP package installation using \textbf{manifest.json}, Docker coordination, inference job control, observability, graphical integration in ELAN, and functional validation.

The results confirmed the technical feasibility of the integration. The middleware installed models, verified their availability, executed inference requests, reported job states, distinguished functional failures from timeouts, and preserved diagnostic metrics. Black-box and white-box tests showed stable behavior in scenarios involving disabled models, missing videos, stopped containers, duplicate installations, and invalid responses. The main contribution is a local, decoupled, and extensible architecture that incorporates AI models into the ELAN workflow while preserving data sovereignty and clearly separating the responsibilities of linguistic annotation, orchestration, and inference.

\textbf{KEYWORDS:} ELAN, middleware, service orchestration, Artificial Intelligence, Ecuadorian Sign Language.
\setcounter{frontmatterlastpage}{\value{page}}